\section{Technical Approach}

The Phase 2 technical approach preserves the Phase 1 architecture while expanding the sensing and evaluation layers.

\subsection{Signal Model}

The baseline signal model remains a narrowband ULA model. For a single source, the array snapshot vector can be written as
\begin{equation}
    \snapshot[m] = \steer(\theta)s[m] + \noise[m],
\end{equation}
where \(\steer(\theta)\) is the ULA steering vector, \(s[m]\) is the source waveform, and \(\noise[m]\) is complex additive noise. The sample covariance matrix is estimated as
\begin{equation}
    \covest = \frac{1}{M}\sum_{m=1}^{M}\snapshot[m]\snapshot^{\Herm}[m].
\end{equation}

This keeps the mathematical foundation aligned with Phase 1 while allowing the experimental conditions to become more difficult.

\subsection{Feature Extraction}

Phase 2 will begin with the Phase 1 beamspace feature representation. Then it will evaluate at least one additional covariance-derived feature set. Candidate features include diagonal/off-diagonal covariance magnitudes, normalized correlation features, or a compact MUSIC-spectrum summary if a classical reference estimator is added.

\subsection{Inference Models}

The continuous nearest-centroid baseline will remain the reference classifier. The spike-based path will use the same conceptual transformation as Phase 1:
\begin{equation}
    \features \rightarrow \spikes[t] \rightarrow \membrane[t].
\end{equation}

For class \(c\), the LIF-inspired temporal state is modeled as
\begin{equation}
    v_c[t] = \leak v_c[t-1] + i_c[t],
\end{equation}
where \(\leak\) is the leak factor and \(i_c[t]\) is the spike-driven class evidence. The final predicted sector is obtained from the class with the largest accumulated membrane state.

\subsection{Embedded-AI Interpretation}

The key embedded-AI interpretation is that spike-based inference should be judged by a multi-objective trade-off, not by accuracy alone. The relevant outputs are:

\begin{itemize}
    \item classification accuracy,
    \item accuracy loss relative to the continuous baseline,
    \item stable-decision latency,
    \item spike count and sparsity,
    \item robustness versus SNR and angular ambiguity.
\end{itemize}
